\documentclass[]{article}

\usepackage{amsmath}
\usepackage{multirow}
\usepackage{natbib}
\usepackage{graphicx} 


\begin{document}

The integration of Artificial Intelligence (AI) into the global economy marks a significant technological inflection point, distinguished from prior automation waves by its capacity to engage with cognitive tasks previously exclusive to human intellect. This advancement introduces a fundamental ambiguity into the workplace, where the potential for augmented productivity and innovation coexists with the specter of skill obsolescence and job displacement. For employees, this transition generates a complex psychological landscape of uncertainty, hope, and anxiety. Consequently, the central challenge for contemporary organizations is not merely to adopt new technologies, but to navigate this transition in a manner that harnesses the benefits of AI while cultivating workforce stability and psychological well-being. Understanding the diverse ways employees perceive AI's impact on their professional future is a critical first step in this process.

Prevailing analyses of AI's effect on job security often rely on variable-centered methods that measure aggregate shifts in employee sentiment. While useful for macroeconomic forecasting, such approaches treat the workforce as a monolithic entity and risk obscuring crucial underlying dynamics. An average measure of job security might remain stable, for instance, while masking a growing polarization between a cohort of employees who feel empowered by AI and another that feels increasingly threatened. This heterogeneity is not random noise but reflects distinct psychological trajectories shaped by personal dispositions, job roles, and, critically, the organizational environment. Without a granular understanding of these divergent response patterns, policies designed to support employees during this transition risk being misaligned, inefficient, or altogether ineffective, failing to address the specific anxieties and aspirations of distinct workforce segments.

This paper addresses this analytical gap by adopting a person-centered approach to model the heterogeneity in employee responses to AI. We move beyond conventional methods to ask whether distinct, unobserved subgroups of employees can be identified based on their perceived job security trajectories. To achieve this, we apply Latent Class Analysis to survey data from a large sample of employees, using the joint distribution of their current and three-year expected job security as indicators. This statistical method allows us to empirically derive and characterize distinct cohorts from the data, revealing a more nuanced typology of workforce adaptation. Our analysis uncovers three primary groups: a "Resiliently Optimistic" class experiencing high and stable security, an "Anxiously Declining" class perceiving a significant threat, and a "Stagnant Neutral" class marked by uncertainty.

Having identified these distinct trajectories, our primary objective is to determine which organizational factors predict an employee's membership in one class over another. We hypothesize that an organization's policies and culture function as active catalysts that can steer employees toward either resilience or anxiety. Using a multinomial logistic regression model, we investigate the influence of specific, actionable organizational levers, including direct employee involvement in AI development, the nature of training programs, non-monetary incentives, and prevailing cultural attributes. By isolating the policies most strongly associated with positive psychological outcomes, this study provides a data-driven framework for organizations seeking to decouple technological advancement from job insecurity. Ultimately, our work aims to shift the discourse from speculating whether AI will replace jobs to understanding how organizations can proactively construct an environment where employees and AI can coexist productively and securely.

\end{document}
                